\documentclass[11pt]{article}
\usepackage[margin=1in]{geometry}
\usepackage{enumitem}
\usepackage[hidelinks]{hyperref}
\setlength{\parskip}{4pt}
\begin{document}
\begin{center}
{\Large\bfseries Response to the Editor}\\[6pt]
Manuscript: \emph{Causal graph recommendation under exposure bias: removing the repeated-walk bias of doubly robust propagation}\\
Journal: Information Processing \& Management
\end{center}

Dear Dr.\ Jansen,

Thank you for reading the manuscript and for the detailed instructions. We have revised the
manuscript and the supplementary material accordingly. Below, each point of your message is
quoted in short form, followed by what we did and where. Section numbers refer to the revised
manuscript; ``Supplement S$n$'' refers to the revised Supplementary Material.

\section*{Summary of the main changes}
\begin{itemize}[leftmargin=*]
\item \textbf{Length.} The manuscript is shortened from about 13{,}000 to 9{,}300 words of text, and from 79
to 53 pages in the review layout (34 pages of text, followed by references, declarations and the
float pages). The proofs, the experimental details and the secondary results (sparse logs, exposure
concentration, stakeholders and regulation) moved to the Supplementary Material.
\item \textbf{New Section 2} ``Research objectives and contributions'' with four objectives and four research questions.
\item \textbf{New Section 7.1} ``Datasets'' with a table describing all four datasets, and Section 7.2 ``Baselines''.
\item \textbf{New Section 9} ``Discussion and implications'' with separate theoretical implications, practical implications and limitations.
\item \textbf{Abstract} rewritten: 216 words, with dataset sizes and numerical results.
\item \textbf{Highlights} rewritten: five items of 19 to 24 words.
\item \textbf{Title} changed, to name the field (causal recommendation) and the contribution.
\end{itemize}

\section*{Point-by-point response}

\paragraph{1. Reviewing for IP\&M.}
Noted. The author will accept review invitations from the journal.

\paragraph{2. ``Highlights are not highlights \dots\ or are too long; about 15 to 30 words each.''}
We rewrote the five highlights so that each states one finding or contribution in 19 to 24 words,
for example the exact bias of debiased graph propagation beyond one hop, the 18--29\% overestimate of
utility that the correction removes on a log with known exposure, and the unchanged accuracy over
four datasets and up to 26 baselines (file \texttt{highlights}). They include the negative result
that the guarantee fails when degree weights come from the log. (Elsevier's general guidance
suggests at most 85 characters per highlight; we followed your length instruction.)

\paragraph{3. ``The abstract does not highlight the specifics \dots\ too much background; add details; shorter is better.''}
The abstract is rewritten from 250 to 216 words. Background is reduced to two sentences. It now
gives the data (four datasets, 290 to 15{,}400 users, up to 10.3 million logged interactions; a
semi-synthetic KuaiRec log), the comparison (up to 26 tuned baselines) and the main numbers: the
three-hop error reduced by a factor of 4 to 12, the 18--29\% utility overestimate removed under oracle
conditions and 100--540\% overestimate with estimated propensities, accuracy differences of at most
0.0014 nDCG, a loss of up to 0.06 nDCG for the sample-split variant, and the 0.004 nDCG margin on KuaiRec
that lies within the seed spread of the best trained model.

\paragraph{4. ``An explicit research objective(s), preferably as a separate section.''}
New \textbf{Section 2, Research objectives and contributions}. It states four objectives (O1
estimator, O2 conditions, O3 value, O4 invariance) with four research questions (RQ1--RQ4), gives
short answers, and lists the contributions. The Results (Section 8) and the Discussion (Section 9)
are organised by the same research questions.

\paragraph{5. ``Questions about your dataset/sample; explicitly describe the dataset(s) in the manuscript.''}
New \textbf{Section 7.1, Datasets}, with Table 2 (users, items, size and density of the biased log,
what the observation indicator means, the unbiased test set, candidates per user and number of
evaluated users) and a paragraph per dataset: Coat (290 users, 300 items, 6{,}960 ratings; 16 random
ratings per user for test), Yahoo!\ R3 (15{,}400 users, 1{,}000 songs, 311{,}704 ratings; 54{,}000
random ratings of 5{,}400 users), KuaiRand-Pure (10{,}000 sampled users, 7{,}583 videos, 511{,}863
impressions; random impressions as test) and KuaiRec (7{,}176 users, 10{,}728 items, 10.3 million views;
a fully observed 1{,}411 $\times$ 3{,}327 matrix as test). We state how positives are defined, how
candidates are chosen, and the metric values of a random ranking. We also state that on Coat and
Yahoo!\ R3 the ``exposure'' is user self-selection, not a platform's display policy.

\paragraph{6. ``No comparison to a state-of-the-art baseline; baselines stale (use or address not using LLMs); unclear whether baselines were run or taken from the literature; same datasets.''}
New \textbf{Section 7.2, Baselines}. (a) \emph{We ran every baseline ourselves}, on the same four
datasets, the same splits and the same validation data; no number is taken from the literature. (b)
The comparison includes recent graph, contrastive and debiased models: BSPM (2023), SimGCL, r-AdjNorm
and NAVIP (2022), MRDR (2021) and GF-CF (2021), together with LightGCN, EASE, iALS, DR-JL, MACR and PDA;
each is tuned per split, and the number of configurations is reported (Supplement S2). (c) A new
paragraph explains the baselines we did not run (StableDR, DICE, CausE, AutoDebias) and why. (d)
\emph{Large language models}: we did not compare with LLM-based recommenders. Our pipeline uses only
interaction data, none of the four benchmarks provides item text together with a random-exposure test
set, and the question studied is the estimation of a propagation operator under exposure bias, for which
the relevant comparators are graph operators. We state this explicitly in Section 7.2 and as a limitation
in Section 9.3. We recognise that this is a genuine limitation and do not claim the method outperforms
LLM-based recommenders. (e) We additionally report new controls for the claim that gains come from the DR
graph: propagation over the imputation alone, equal-budget comparisons and seed variation (Sections 8.3, 9.3).

\paragraph{7. ``Highlight the theoretical and practical implications; how does this research differ from existing work; explicit discussion section.''}
New \textbf{Section 9, Discussion and implications}, with Section 9.1 theoretical implications
(propagation-level debiasing must consider the operator; a guarantee for fixed nuisances is a statement
about the whole pipeline, including degree weights; unbiasedness for the full-exposure target does not imply
better rankings), Section 9.2 practical implications (when the correction is and is not worth using, what
needs true propensities and nuisances independent of the log, what the audit can and cannot show), and Section
9.3 limitations. The difference from existing work is stated in Section 3 (``Debiasing inside graph
propagation'' and ``Unbiased estimation of polynomials under missingness'') and in Section 2: the reduction itself is classical;
what is new is its exact form for inverse-propensity and DR graph propagation, the exact bias of the
uncorrected operators, and the finding that the guarantees fail under data-dependent degree weights.

\paragraph{8. ``Thorough copy editing; many readers are not native English speakers.''}
We rewrote the introduction, related work, experimental design, results, discussion and conclusion in plain
language with shorter sentences, defined the central symbols in the introduction, and edited the prose of the method section
for clarity without changing any mathematics, citation, number or result (verified by an automated check of
formulas, citations, labels and numbers before and after). We would welcome a further language check at
the proof stage.

\paragraph{9. ``The manuscript is WAY too long; shorten it.''}
See the summary: 13{,}000 $\to$ 9{,}300 words of text and 79 $\to$ 53 pages in the review layout, with 34 pages of
text. Concretely: the related-work section is about 40\% shorter (1{,}320 to 770 words); the simulation, accuracy and audit
sections are condensed to the findings needed for the four research questions; the variance and clipping theorem
with its proof, the experimental details and the proofs of all theorems moved to the Supplementary Material (S1--S3);
and the sparse-log experiment, the exposure-concentration and clip analyses, and the stakeholder and regulation discussion
moved to Supplement S5.

\paragraph{10. ``The reference list needs tidying; missing items or formatting issues; use a standard style (APA).''}
We audited every cited entry for author, title, venue, year, volume, pages and DOI, completed the missing
pages, fixed capitalisation of titles (e.g.\ M\"obius, ``A worrying analysis''), and the printed list contains only the 65 works cited in the
manuscript; we use the Elsevier author--year (Harvard) style with
consistent formatting throughout. This style is close to, but not identical to, APA 7; if the journal requires strict
APA formatting we will convert at the revision stage.

\paragraph{11. ``Review recent IP\&M articles on similar topics.''}
The related work already cites recent IP\&M articles on debiasing, fairness and privacy (for example on
popularity debiasing from user and item perspectives, controllable popularity bias, invariant debiasing with
biased imputation, causal deconfounding of multi-criteria ratings, fairness measurement, provider fairness and
differential privacy for matrix factorisation), and we kept them. We restructured the paper along the usual
research-article layout (introduction, objectives, related work, method, experimental design, results, discussion
and implications, conclusion) and moved supporting material to the supplement.

\paragraph{12. Author guidelines.}
The abstract is under 250 words, the highlights follow your length instruction, and the title page contains the
declarations of funding, competing interests, generative AI use, CRediT roles, ethics and data availability, with the
public repository for code and result files.

\section*{Other changes made during the revision}
\begin{itemize}[leftmargin=*]
\item New tests: equivalence (TOST) for the key comparisons, a Nadeau--Bengio test over splits adjusted over the same family as the user-level tests, and a note that identical rankings (Yahoo!\ R3, KuaiRand) cannot be tested (Supplement S5, Table S19).
\item A simulation control that fixes the degree weights a priori shows that cross-fitting alone preserves the correction and that data-dependent degree weights break it (Section 8.1, Table S18).
\item Table 1 now separates DRUP-split (assumptions hold) from cross-fitted DRUP (as run, no guarantee).
\item The claim that the correction ``changes no ranking'' is replaced by ``no detectable change in accuracy'', and the exposure-invariance audit is described as having limited power.
\end{itemize}

\section*{Points we did not address, and why}
\begin{itemize}[leftmargin=*]
\item \textbf{LLM-based and newer baselines.} Not run (Section 7.2). A comparison needs a dataset with item text and a random-exposure test set.
\item \textbf{Strict APA formatting.} The reference list uses the Elsevier author--year style (see point 10).
\item \textbf{Main-text length.} At 34 pages in the review layout the text is still long for some readers. Further shortening would require removing results that answer the research questions.
\end{itemize}

We hope the revised manuscript addresses your concerns, and we are grateful for the opportunity to resubmit.

Sincerely,\\
Quang-Vinh Dang\\
British University Vietnam, Hung Yen, Vietnam
\end{document}
